\uuid{p9K4}
\titre{Espaces vectoriels engendrés dans $\mathbb{R}^3$}

\niveau{} 				%L1, L2, L3, MPSI, MP, PCSI, PC, PSI...
\module{ Espaces vectoriels } 	%Analyse, Algèbre...
\chapitre{Dimension et bases}   			%Continuité, Groupes, Fonctions de plusieurs variables...
\sousChapitre{}			%Optimisation, Diagonalisation d'une matrice, Calcul de dérivées partielles...

\theme{}				%Fonction de répartition, Division euclidienne de polynômes, ...
\auteur{Quentin Liard}
\datecreate{ 2026-02-19}
\organisation{AMSCC}			%AMSCC, Exo7, ...
\difficulte{}			%1, 2, 3, 4 ou 5

\contenu{

\texte{ 
Dans $\mathbb{R}^3$, on considère
$u=(1,1,0)$, $v=(0,1,1)$, $w=(1,2,1)$, $z=(1,0,-1)$.
On pose $F=\mathrm{Vect}(u,v,w)$ et $G=\mathrm{Vect}(u,z)$.
}

\begin{enumerate}
\item   \question{La famille $(u,v,w)$ est-elle libre ?}
\indication{}
\reponse{Non, car $u+v=(1,1,0)+(0,1,1)=(1,2,1)=w$. Donc $w-u-v=0$ : dépendance linéaire.}

\item   \question{Déterminer $\dim F$ et  $\dim G$.}
\indication{}
\reponse{Puisque $w=u+v$, on a $F=\mathrm{Vect}(u,v,w)=\mathrm{Vect}(u,v)$.
Or $u$ et $v$ ne sont pas colinéaires, donc ils sont libres. Ainsi $\dim F=2$.
$G=\mathrm{Vect}(u,z)$. Les vecteurs $u=(1,1,0)$ et $z=(1,0,-1)$ ne sont pas colinéaires, donc ils sont libres. Ainsi $\dim G=2$.
}

%\item   \question{Déterminer.}
%\indication{}
%\reponse{$G=\mathrm{Vect}(u,z)$. Les vecteurs $u=(1,1,0)$ et $z=(1,0,-1)$ ne sont pas colinéaires, donc ils sont libres. Ainsi $\dim G=2$.}

\item   \question{Déterminer $F\cap G$.}
\indication{}
\reponse{Montrons que $F\subset G$ : il suffit de vérifier que $v\in G$ (car $F=\mathrm{Vect}(u,v)$).
Cherchons $c,d$ tels que $v=c\,u+d\,z$ :
$c(1,1,0)+d(1,0,-1)=(c+d,\ c,\ -d)=(0,1,1)$.
On obtient $c=1$, $-d=1\Rightarrow d=-1$, et $c+d=0$.
Donc $v\in G$ et ainsi $F=\mathrm{Vect}(u,v)\subset \mathrm{Vect}(u,z)=G$.
Comme $\dim F=\dim G=2$, on a en fait $F=G$ et donc $F\cap G=F$ (et $\dim(F\cap G)=2$).}

\item   \question{Déterminer $\dim(F+G)$.}
\indication{}
\reponse{Comme $F=G$, on a $F+G=F$, donc $\dim(F+G)=\dim F=2$.}

\end{enumerate}

}